% !TEX spellcheck = en_US

% ~~~~~~~~~~~~~~~~~~~~~~~~~~~~~~~~~~~~~~~ V
% (NC) 2026 Haluk Bingol
% github.com/halukbingol/LaTeX-Templates
%
% Licensees may copy, distribute, display, and perform the work and make derivative works
% and remixes based on it only for non-commercial purposes.
% ~~~~~~~~~~~~~~~~~~~~~~~~~~~~~~~~~~~~~~~ A

% v20261107T1200   user manual of hblibCodeoutput.sty v20261104T2350
%
% hblibCodeoutput-UserManual.tex -- every feature of hblibCodeoutput.sty, as a deck:
% the LaTeX source next to the result.
% Put it next to hblibPresentation.sty and hblibCodeoutput.sty, then compile twice:
%   pdflatex hblibCodeoutput-UserManual
%   pdflatex hblibCodeoutput-UserManual
% It writes the example files Hello.java, Hello.txt, Grades.java, Grades.txt and Session.txt
% next to itself (see "Example files" below). The outputs are real: the programs were compiled
% and run with JDK 21.

\documentclass[aspectratio=169,10pt,compress]{beamer}




% ~~~~~~~~~~~~~~~~~~~~~~~~~~~~~~~~~~~~~~~~~~~~~~~~~~~~~~~~~~~~~~~~~ V hblib
	\usepackage{hblibPresentation}     % in a deck: \usepackage{../../hblib/hblibPresentation}
	\usepackage{hblibCodeoutput}       % in a deck: \usepackage{../../hblib/hblibCodeoutput}
% ~~~~~~~~~~~~~~~~~~~~~~~~~~~~~~~~~~~~~~~~~~~~~~~~~~~~~~~~~~~~~~~~~ A hblib




% +++++++++++++++++++++++++++++++++++++++ V specific
% texcode: LaTeX source in a light box (this manual only)
\lstnewenvironment{texcode}{\lstset{style=texsnippet, style=utf8letters, numbers=none, columns=fullflexible,
	keepspaces=true, breaklines=true, texcsstyle=*\color{blue!55!black}, commentstyle=\color{gray}\itshape,
	aboveskip=2pt, belowskip=2pt}}{}
% +++++++++++++++++++++++++++++++++++++++ A specific




% ~~~~~~~~~~~~~~~~~~~~~~~~~~~~~~~~~~~~~~~ V example files
% real programs and their outputs (compiled and run with JDK 21); in a deck they live in codeoutput/
\begin{filecontents*}[overwrite]{Hello.java}
public class Hello {
    public static void main(String[] args) {
        String name = "Ayşe";
        System.out.println("Merhaba, " + name + "!");
        System.out.println("Price: 25 ₺");
    }
}
\end{filecontents*}
\begin{filecontents*}[overwrite]{Hello.txt}
Merhaba, Ayşe!
Price: 25 ₺
\end{filecontents*}
\begin{filecontents*}[overwrite]{Grades.java}
public class Grades {
    public static void main(String[] args) {
        int[] marks = {72, 85, 90, 64, 78};
        int sum = 0;
        for (int m : marks) {
            sum += m;
        }
        double average = (double) sum / marks.length;
        System.out.println("count   = " + marks.length);
        System.out.println("sum     = " + sum);
        System.out.println("average = " + average);
    }
}
\end{filecontents*}
\begin{filecontents*}[overwrite]{Grades.txt}
count   = 5
sum     = 389
average = 77.8
\end{filecontents*}
\begin{filecontents*}[overwrite]{Session.txt}
$ javac Hello.java
$ ls
Hello.class
Hello.java
$ java Hello
Merhaba, Ayşe!
Price: 25 ₺
\end{filecontents*}
% ~~~~~~~~~~~~~~~~~~~~~~~~~~~~~~~~~~~~~~~ A example files




\title[hblibCodeoutput: User Manual]{
	hblibCodeoutput: User Manual
}
\subtitle{
	for hblibCodeoutput.sty v20261104T2350
}
\author[Bingol]{Haluk Bingol}

\institute[FCIS]{
	Faculty of Computer and Information Sciences\\
	Yeditepe University
}
\date{{\tiny \hbTimeStamp}}

\begin{document}




% ~~~~~~~~~~~~~~~~~~~~~~~~~~~~~~~~~~~~~~ V frm
\begin{frame}[t,fragile]
	\titlepage
\end{frame}
% ~~~~~~~~~~~~~~~~~~~~~~~~~~~~~~~~~~~~~~ A frm




% ~~~~~~~~~~~~~~~~~~~~~~~~~~~~~~~~~~~~~~ V frm
\begin{frame}[t,fragile] \frametitle{Outline}
	\begin{columns}[T]
		\begin{column}{0.48\textwidth}
			\tableofcontents[sections={1-4}, hideallsubsections]
		\end{column}
		\begin{column}{0.48\textwidth}
			\tableofcontents[sections={5-7}, hideallsubsections]
		\end{column}
	\end{columns}
\end{frame}
% ~~~~~~~~~~~~~~~~~~~~~~~~~~~~~~~~~~~~~~ A frm




% ~~~~~~~~~~~~~~~~~~~~~~~~~~~~~~~~~~~~~~ V
\hSection[Start]{Getting started}{Files, paths, real outputs}
% ~~~~~~~~~~~~~~~~~~~~~~~~~~~~~~~~~~~~~~ A




% ~~~~~~~~~~~~~~~~~~~~~~~~~~~~~~~~~~~~~~ V frm
\begin{frame}[t,fragile] \frametitle{The idea: code and output stay in files}
	\begin{columns}[T]
		\begin{column}{0.50\textwidth}

% +++++++++++++++++++++++++++++++++++++++ V
%: -Code.
\begin{texcode}
\usepackage{../../hblib/hblibPresentation}
\usepackage{../../hblib/hblibCodeoutput}

\begin{frame}[t,fragile] \frametitle{Sum}
  \lstcoderightstep{codeoutput/Grades.java}{Java}
                   {codeoutput/Grades.txt}
\end{texcode}
% +++++++++++++++++++++++++++++++++++++++ A

		\end{column}
		\begin{column}{0.46\textwidth}
			\small A lecture folder:

% +++++++++++++++++++++++++++++++++++++++ V
%: -Code.
\begin{lstlisting}[style=outputstyle, frame=single, rulecolor=\color{gray!40}]
lecture/
  cis-array1D.tex
  run_examples.sh
  codeoutput/
    Grades.java    the program
    Grades.in      keyboard input
    Grades.txt     what it printed
\end{lstlisting}
% +++++++++++++++++++++++++++++++++++++++ A

		\end{column}
	\end{columns}

	\begin{itemize}\small
		\item Every macro takes \hEmph{relative paths}, relative to the .tex file; nothing is copied into the .tex
		\item Frames with these macros must be \hCode{[fragile]}
		\item Load hblibCodeoutput after hblibPresentation; it is made for beamer
		(outside beamer it warns, and the \hCode{*step} macros do not work)
	\end{itemize}
\end{frame}
% ~~~~~~~~~~~~~~~~~~~~~~~~~~~~~~~~~~~~~~ A frm




% ~~~~~~~~~~~~~~~~~~~~~~~~~~~~~~~~~~~~~~ V frm
\begin{frame}[t,fragile] \frametitle{Real outputs: run the programs, keep what they print}
	\small Outputs are never typed by hand. A script runs every program and writes its \hCode{.txt}:

% +++++++++++++++++++++++++++++++++++++++ V
%: -Code.
\begin{lstlisting}[style=codesnippet, language=bash]
#!/bin/bash
# run_examples.sh: compile and run every program in codeoutput/, keep its output
cd codeoutput
for f in *.java; do
    name="${f%.java}"
    javac "$f"
    if [ -f "$name.in" ]; then java "$name" < "$name.in" > "$name.txt"
    else                       java "$name"              > "$name.txt"; fi
done
\end{lstlisting}
% +++++++++++++++++++++++++++++++++++++++ A

	\begin{itemize}\small
		\item Change a program, run \hCode{sh run\_examples.sh}, compile the deck: code and output always match
		\item A missing file gives a red box, not an error (see ``Notes''), so a deck compiles before its outputs exist
	\end{itemize}
\end{frame}
% ~~~~~~~~~~~~~~~~~~~~~~~~~~~~~~~~~~~~~~ A frm




% ~~~~~~~~~~~~~~~~~~~~~~~~~~~~~~~~~~~~~~ V
\hSection[Blocks]{Building blocks}{\hCode{\textbackslash lstcode}, \hCode{\textbackslash lstoutput}, \hCode{\textbackslash lstterminal}}
% ~~~~~~~~~~~~~~~~~~~~~~~~~~~~~~~~~~~~~~ A




% ~~~~~~~~~~~~~~~~~~~~~~~~~~~~~~~~~~~~~~ V frm
\begin{frame}[t,fragile] \frametitle{\hCode{\textbackslash lstcode[options]\{code path\}\{Language\}}}
	\begin{columns}[T]
		\begin{column}{0.46\textwidth}

% +++++++++++++++++++++++++++++++++++++++ V
%: -Code.
\begin{texcode}
\lstcode{Grades.java}{Java}
\end{texcode}
% +++++++++++++++++++++++++++++++++++++++ A

			\footnotesize
			\begin{itemize}
				\item Title bar: the path, and the language on the right
				\item Line numbers on the left
				\item \hCode{Language} is a listings language: \hCode{Java}, \hCode{C}, \hCode{Python}, \hCode{bash}, \dots
			\end{itemize}
		\end{column}
		\begin{column}{0.50\textwidth}
			\lstcode[basicstyle=\ttfamily\scriptsize]{Grades.java}{Java}
		\end{column}
	\end{columns}
\end{frame}
% ~~~~~~~~~~~~~~~~~~~~~~~~~~~~~~~~~~~~~~ A frm




% ~~~~~~~~~~~~~~~~~~~~~~~~~~~~~~~~~~~~~~ V frm
\begin{frame}[t,fragile] \frametitle{Part of a file: \hCode{firstline}, \hCode{lastline}, \hCode{firstnumber}}

% +++++++++++++++++++++++++++++++++++++++ V
%: -Code.
\begin{texcode}
\lstcode[firstline=3, lastline=7]{Grades.java}{Java}                  % numbered 1 ... 5
\lstcode[firstline=3, lastline=7, firstnumber=3]{Grades.java}{Java}   % numbered 3 ... 7, as in the file
\end{texcode}
% +++++++++++++++++++++++++++++++++++++++ A

	\begin{columns}[T]
		\begin{column}{0.48\textwidth}
			\lstcode[basicstyle=\ttfamily\scriptsize, firstline=3, lastline=7]{Grades.java}{Java}
		\end{column}
		\begin{column}{0.48\textwidth}
			\lstcode[basicstyle=\ttfamily\scriptsize, firstline=3, lastline=7, firstnumber=3]{Grades.java}{Java}
		\end{column}
	\end{columns}

	\small With \hCode{firstnumber}, ``line 5'' in the text is line 5 of the file.
	Other options: \hCode{basicstyle=\textbackslash ttfamily\textbackslash scriptsize},
	\hCode{linerange=\{1-2,9-11\}}, \hCode{numbers=none}.
\end{frame}
% ~~~~~~~~~~~~~~~~~~~~~~~~~~~~~~~~~~~~~~ A frm




% ~~~~~~~~~~~~~~~~~~~~~~~~~~~~~~~~~~~~~~ V frm
\begin{frame}[t,fragile] \frametitle{\hCode{\textbackslash lstoutput} and \hCode{\textbackslash lstterminal}}
	\begin{columns}[T]
		\begin{column}{0.48\textwidth}

% +++++++++++++++++++++++++++++++++++++++ V
%: -Code.
\begin{texcode}
\lstoutput{Grades.txt}       % output only
\end{texcode}
% +++++++++++++++++++++++++++++++++++++++ A

			\lstoutput{Grades.txt}

			\footnotesize A dark box, like a terminal: ``Output'' and the path in the title bar.
		\end{column}
		\begin{column}{0.48\textwidth}

% +++++++++++++++++++++++++++++++++++++++ V
%: -Code.
\begin{texcode}
\lstterminal{Session.txt}    % a terminal session
\end{texcode}
% +++++++++++++++++++++++++++++++++++++++ A

			\lstterminal{Session.txt}

			\footnotesize Lines starting with \hCode{\$ } are the commands; they are highlighted.
		\end{column}
	\end{columns}
\end{frame}
% ~~~~~~~~~~~~~~~~~~~~~~~~~~~~~~~~~~~~~~ A frm




% ~~~~~~~~~~~~~~~~~~~~~~~~~~~~~~~~~~~~~~ V
\hSection[One slide]{Code and output on one slide}{\hCode{\textbackslash lstcodedown}, \hCode{\textbackslash lstcoderight}}
% ~~~~~~~~~~~~~~~~~~~~~~~~~~~~~~~~~~~~~~ A




% ~~~~~~~~~~~~~~~~~~~~~~~~~~~~~~~~~~~~~~ V frm
\begin{frame}[t,fragile] \frametitle{\hCode{\textbackslash lstcodedown}: output below}

% +++++++++++++++++++++++++++++++++++++++ V
%: -Code.
\begin{texcode}
\lstcodedown[code options][output options]{code path}{Language}{output path}
\lstcodedown[basicstyle=\ttfamily\tiny, firstline=8, lastline=11, firstnumber=8]{Grades.java}{Java}{Grades.txt}
\end{texcode}
% +++++++++++++++++++++++++++++++++++++++ A

	\lstcodedown[basicstyle=\ttfamily\tiny, firstline=8, lastline=11, firstnumber=8]{Grades.java}{Java}{Grades.txt}
\end{frame}
% ~~~~~~~~~~~~~~~~~~~~~~~~~~~~~~~~~~~~~~ A frm




% ~~~~~~~~~~~~~~~~~~~~~~~~~~~~~~~~~~~~~~ V frm
\begin{frame}[t,fragile] \frametitle{\hCode{\textbackslash lstcoderight}: output on the right}

% +++++++++++++++++++++++++++++++++++++++ V
%: -Code.
\begin{texcode}
\lstcoderight[code width][code options][output options]{code path}{Language}{output path}
\lstcoderight[0.66][basicstyle=\ttfamily\scriptsize]{Grades.java}{Java}{Grades.txt}
\end{texcode}
% +++++++++++++++++++++++++++++++++++++++ A

	\lstcoderight[0.66][basicstyle=\ttfamily\scriptsize]{Grades.java}{Java}{Grades.txt}

	\small \hCode{code width}: a fraction of \hCode{\textbackslash textwidth} (default 0.56); the output gets the rest.
\end{frame}
% ~~~~~~~~~~~~~~~~~~~~~~~~~~~~~~~~~~~~~~ A frm




% ~~~~~~~~~~~~~~~~~~~~~~~~~~~~~~~~~~~~~~ V
\hSection[Two slides]{Output on the second slide}{\hCode{\textbackslash lstcodedownstep}, \hCode{\textbackslash lstcoderightstep}, \hCode{\textbackslash lstcodeoverlaystep}}
% ~~~~~~~~~~~~~~~~~~~~~~~~~~~~~~~~~~~~~~ A




% ~~~~~~~~~~~~~~~~~~~~~~~~~~~~~~~~~~~~~~ V frm
\begin{frame}[t,fragile] \frametitle{\hCode{\textbackslash lstcoderightstep}: ask first, then show}

% +++++++++++++++++++++++++++++++++++++++ V
%: -Code.
\begin{texcode}
\lstcoderightstep[0.66][basicstyle=\ttfamily\scriptsize]{Grades.java}{Java}{Grades.txt}
\small \hQuestion{What is the average?}  \uncover<2->{$389 / 5 = 77.8$}
\end{texcode}
% +++++++++++++++++++++++++++++++++++++++ A

	\lstcoderightstep[0.66][basicstyle=\ttfamily\scriptsize]{Grades.java}{Java}{Grades.txt}

	\small \hQuestion{What is the average?} \uncover<2->{$389 / 5 = 77.8$}

	\footnotesize Slide 1: the code. Slide 2: the output too.
	Its space is kept on slide 1, so nothing jumps; the rest of the frame is not delayed.
\end{frame}
% ~~~~~~~~~~~~~~~~~~~~~~~~~~~~~~~~~~~~~~ A frm




% ~~~~~~~~~~~~~~~~~~~~~~~~~~~~~~~~~~~~~~ V frm
\begin{frame}[t,fragile] \frametitle{\hCode{\textbackslash lstcodedownstep}: output below, on slide 2}

% +++++++++++++++++++++++++++++++++++++++ V
%: -Code.
\begin{texcode}
\lstcodedownstep[basicstyle=\ttfamily\scriptsize, firstline=3, lastline=8, firstnumber=3]{Grades.java}{Java}{Grades.txt}
\end{texcode}
% +++++++++++++++++++++++++++++++++++++++ A

	\lstcodedownstep[basicstyle=\ttfamily\scriptsize, firstline=3, lastline=8, firstnumber=3]{Grades.java}{Java}{Grades.txt}
\end{frame}
% ~~~~~~~~~~~~~~~~~~~~~~~~~~~~~~~~~~~~~~ A frm




% ~~~~~~~~~~~~~~~~~~~~~~~~~~~~~~~~~~~~~~ V frm
\begin{frame}[t,fragile] \frametitle{\hCode{\textbackslash lstcodeoverlaystep}: the output floats on top}

% +++++++++++++++++++++++++++++++++++++++ V
%: -Code.
\begin{texcode}
\lstcodeoverlaystep[code options][output options][output width]{x,y}{code path}{Language}{output path}
\lstcodeoverlaystep[basicstyle=\ttfamily\scriptsize][][5cm]{9cm,3.6cm}{Grades.java}{Java}{Grades.txt}
\end{texcode}
% +++++++++++++++++++++++++++++++++++++++ A

	\lstcodeoverlaystep[basicstyle=\ttfamily\scriptsize][][5cm]{9cm,3.6cm}{Grades.java}{Java}{Grades.txt}

	\footnotesize \hCode{x,y}: the output's upper-left corner, from the slide's upper-left corner (x right, y down).
	Output width default \hCode{0.45\textbackslash textwidth}. Use it when the code fills the frame.
\end{frame}
% ~~~~~~~~~~~~~~~~~~~~~~~~~~~~~~~~~~~~~~ A frm




% ~~~~~~~~~~~~~~~~~~~~~~~~~~~~~~~~~~~~~~ V frm
\begin{frame}[t,fragile] \frametitle{\hCode{\textbackslash slidegrid}: finding x and y}
	\slidegrid

% +++++++++++++++++++++++++++++++++++++++ V
%: -Code.
\begin{texcode}
\begin{frame}[t,fragile]
  \slidegrid         % a labeled 1 cm grid over the slide; remove it when done
  ...
\end{texcode}
% +++++++++++++++++++++++++++++++++++++++ A

	\small Read the corner you want from the grid, e.g.\ \hCode{\{9cm,3.6cm\}}.
	It works for any slide size (16\,$\times$\,9\,cm here, 12.8\,$\times$\,9.6\,cm for 4:3).
\end{frame}
% ~~~~~~~~~~~~~~~~~~~~~~~~~~~~~~~~~~~~~~ A frm




% ~~~~~~~~~~~~~~~~~~~~~~~~~~~~~~~~~~~~~~ V
\hSection[Typed]{Typed and inline code}{codesnippet, texsnippet, terminalstyle; \hCode{\textbackslash hCodeJ}, \hCode{\textbackslash hCodeB}}
% ~~~~~~~~~~~~~~~~~~~~~~~~~~~~~~~~~~~~~~ A




% ~~~~~~~~~~~~~~~~~~~~~~~~~~~~~~~~~~~~~~ V frm
\begin{frame}[t,fragile] \frametitle{Short code typed in the .tex}
	\begin{columns}[T]
		\begin{column}{0.48\textwidth}

% +++++++++++++++++++++++++++++++++++++++ V
%: -Code.
\begin{texcode}
\begin{lstlisting}[style=codesnippet, language=Java]
double[] t = {14.0, 16.5};   // one name
System.out.println(t[1]);    // index 1
\end{lstlisting}
\end{texcode}
% +++++++++++++++++++++++++++++++++++++++ A

% +++++++++++++++++++++++++++++++++++++++ V
%: -Code.
\begin{lstlisting}[style=codesnippet, language=Java]
double[] t = {14.0, 16.5};   // one name
System.out.println(t[1]);    // index 1
\end{lstlisting}
% +++++++++++++++++++++++++++++++++++++++ A

		\end{column}
		\begin{column}{0.48\textwidth}

% +++++++++++++++++++++++++++++++++++++++ V
%: -Code.
\begin{texcode}
\begin{lstlisting}[style=terminalstyle,
  backgroundcolor=\color{termbg},
  basicstyle=\ttfamily\footnotesize\color{termfg}]
$ java Hello
Merhaba, Ayşe!
\end{lstlisting}
\end{texcode}
% +++++++++++++++++++++++++++++++++++++++ A

% +++++++++++++++++++++++++++++++++++++++ V
%: -Code.
\begin{lstlisting}[style=terminalstyle, backgroundcolor=\color{termbg}, basicstyle=\ttfamily\footnotesize\color{termfg}]
$ java Hello
Merhaba, Ayşe!
\end{lstlisting}
% +++++++++++++++++++++++++++++++++++++++ A

		\end{column}
	\end{columns}

	\small Styles: \hCode{codesnippet} (framed, light), \hCode{codestyle} (numbered, no frame),
	\hCode{outputstyle}, \hCode{terminalstyle}, \hCode{texsnippet} (LaTeX source, as in the boxes of this manual).
\end{frame}
% ~~~~~~~~~~~~~~~~~~~~~~~~~~~~~~~~~~~~~~ A frm




% ~~~~~~~~~~~~~~~~~~~~~~~~~~~~~~~~~~~~~~ V frm
\begin{frame}[t,fragile] \frametitle{Inline code}

% +++++++++++++++++++++++++++++++++++++++ V
%: -Code.
\begin{texcode}
\hCodeJ{int n = a.length;}    \hCodeB{ls -l *.java}    \code{for (int m : marks)}    \hCode{int n = a.length;}
\end{texcode}
% +++++++++++++++++++++++++++++++++++++++ A

	\begin{center}
		\hCodeJ{int n = a.length;} \quad \hCodeB{ls -l *.java} \quad \code{for (int m : marks)} \quad \hCode{int n = a.length;}
	\end{center}

% +++++++++++++++++++++++++++++++++++++++ V tbl
%: -Tbl. inline
\begin{table}
		\footnotesize
		\begin{tabular}{@{}L{0.20\textwidth}L{0.36\textwidth}L{0.36\textwidth}@{}}
			\toprule
			\textbf{Command} & \textbf{Gives} & \textbf{Where} \\
			\midrule
			\hCode{\textbackslash hCodeJ}, \hCode{\textbackslash code} & Java, syntax colored, typed as is & text of a \hCode{[fragile]} frame \\
			\hCode{\textbackslash hCodeB} & bash, typed as is & text of a \hCode{[fragile]} frame \\
			\hCode{\textbackslash hCode} (hblibPresentation) & plain, colored; escape \hCode{\textbackslash \_ \textbackslash \{ \textbackslash \}} & everywhere, also frame titles \\
			\bottomrule
		\end{tabular}
\end{table}%
% +++++++++++++++++++++++++++++++++++++++ A tbl

	\small \hCode{\textbackslash hCodeJ} and \hCode{\textbackslash hCodeB} cannot be used in frame titles,
	and their text must not contain \hCode{\%} or \hCode{\#}.
\end{frame}
% ~~~~~~~~~~~~~~~~~~~~~~~~~~~~~~~~~~~~~~ A frm




% ~~~~~~~~~~~~~~~~~~~~~~~~~~~~~~~~~~~~~~ V
\hSection[Characters]{Characters and colors}{Turkish letters, \hCode{\textbackslash hbLira}, color names}
% ~~~~~~~~~~~~~~~~~~~~~~~~~~~~~~~~~~~~~~ A




% ~~~~~~~~~~~~~~~~~~~~~~~~~~~~~~~~~~~~~~ V frm
\begin{frame}[t,fragile] \frametitle{Turkish letters and currency signs}

% +++++++++++++++++++++++++++++++++++++++ V
%: -Code.
\begin{texcode}
\lstcoderight[0.6][basicstyle=\ttfamily\scriptsize]{Hello.java}{Java}{Hello.txt}
In text: 25\,\hbLira{} and 25\,\hbRupee{}
\end{texcode}
% +++++++++++++++++++++++++++++++++++++++ A

	\lstcoderight[0.6][basicstyle=\ttfamily\scriptsize]{Hello.java}{Java}{Hello.txt}

	\small In text: 25\,\hbLira{} and 25\,\hbRupee{}

	\begin{itemize}\footnotesize
		\item Code and output print Turkish and Latin-1 letters, \texteuro{} £ ¥ °, quotes, dashes and no-break spaces
		\item \hbLira{} and \hbRupee{} are drawn (few TeX fonts have them): \hCode{\textbackslash hbLira}, \hCode{\textbackslash hbRupee}
		\item Wrongly decoded text (``mojibake'', e.g.\ AÄŸaç for Ağaç) also prints, so it can be shown
	\end{itemize}
\end{frame}
% ~~~~~~~~~~~~~~~~~~~~~~~~~~~~~~~~~~~~~~ A frm




% ~~~~~~~~~~~~~~~~~~~~~~~~~~~~~~~~~~~~~~ V frm
\begin{frame}[t,fragile] \frametitle{Colors}
	\begin{columns}[T]
		\begin{column}{0.50\textwidth}

% +++++++++++++++++++++++++++++++++++++++ V tbl
%: -Tbl. colors
\begin{table}
		\scriptsize
		\begin{tabular}{@{}L{0.28\textwidth}L{0.22\textwidth}L{0.36\textwidth}@{}}
			\toprule
			\textbf{Name} & \textbf{RGB} & \textbf{Used for} \\
			\midrule
			\textcolor{codeframe}{\rule{6pt}{6pt}} codeframe & 40, 90, 160    & code box frame, title bar \\
			\fcolorbox{gray}{codebg}{\rule{0pt}{4pt}} codebg & 248, 248, 248 & code box background \\
			\textcolor{codekw}{\rule{6pt}{6pt}} codekw       & 0, 70, 160     & keywords \\
			\textcolor{codestr}{\rule{6pt}{6pt}} codestr     & 30, 130, 60    & strings \\
			\textcolor{codecmt}{\rule{6pt}{6pt}} codecmt     & 120, 120, 120  & comments, line numbers \\
			\textcolor{termbg}{\rule{6pt}{6pt}} termbg       & 35, 38, 45     & output box background \\
			\textcolor{termframe}{\rule{6pt}{6pt}} termframe & 70, 75, 85     & output box frame \\
			\fcolorbox{gray}{termfg}{\rule{0pt}{4pt}} termfg & 225, 230, 235  & output text \\
			\textcolor{termcmd}{\rule{6pt}{6pt}} termcmd     & 255, 214, 102  & terminal commands \\
			\bottomrule
		\end{tabular}
\end{table}%
% +++++++++++++++++++++++++++++++++++++++ A tbl

		\end{column}
		\begin{column}{0.46\textwidth}

% +++++++++++++++++++++++++++++++++++++++ V
%: -Code.
\begin{texcode}
\definecolor{codeframe}{RGB}{150,40,40}
\lstcode{Grades.java}{Java}
\end{texcode}
% +++++++++++++++++++++++++++++++++++++++ A

			{\definecolor{codeframe}{RGB}{150,40,40}
			\lstcode[basicstyle=\ttfamily\tiny, firstline=3, lastline=4, firstnumber=3]{Grades.java}{Java}}

			\footnotesize In the preamble: the whole deck. In a frame: that frame only.
		\end{column}
	\end{columns}
\end{frame}
% ~~~~~~~~~~~~~~~~~~~~~~~~~~~~~~~~~~~~~~ A frm




% ~~~~~~~~~~~~~~~~~~~~~~~~~~~~~~~~~~~~~~ V
\hSection[Notes]{Notes and troubleshooting}\label{sec:notes}
% ~~~~~~~~~~~~~~~~~~~~~~~~~~~~~~~~~~~~~~ A




% ~~~~~~~~~~~~~~~~~~~~~~~~~~~~~~~~~~~~~~ V frm
\begin{frame}[t,fragile] \frametitle{A missing file}

% +++++++++++++++++++++++++++++++++++++++ V
%: -Code.
\begin{texcode}
\lstcoderight{codeoutput/NotYet.java}{Java}{codeoutput/NotYet.txt}
\end{texcode}
% +++++++++++++++++++++++++++++++++++++++ A

	\lstcoderight{codeoutput/NotYet.java}{Java}{codeoutput/NotYet.txt}

	\small A red box names the missing file; the deck still compiles.
	Create the file, or check the path (relative to the .tex file, letters in the right case).
\end{frame}
% ~~~~~~~~~~~~~~~~~~~~~~~~~~~~~~~~~~~~~~ A frm




% ~~~~~~~~~~~~~~~~~~~~~~~~~~~~~~~~~~~~~~ V frm
\begin{frame}[t,fragile] \frametitle{Troubleshooting}

% +++++++++++++++++++++++++++++++++++++++ V tbl
%: -Tbl. troubleshooting
\begin{table}
		\scriptsize
		\begin{tabular}{@{}L{0.38\textwidth}L{0.27\textwidth}L{0.27\textwidth}@{}}
			\toprule
			\textbf{Message or problem} & \textbf{Cause} & \textbf{Fix} \\
			\midrule
			Undefined control sequence \textbackslash tcbuselibrary & a library older than v20261104T2350 & update hblibCodeoutput.sty \\
			errors inside a frame with code & frame not \hCode{fragile} & \hCode{[t,fragile]} \\
			red box ``missing file'' & wrong path, or no output yet & fix the path; run the programs \\
			the output is shown on slide 1 & a one-slide macro & use a \hCode{*step} macro \\
			line numbers start at 1 with \hCode{firstline} & listings counts printed lines & add \hCode{firstnumber=} \\
			the output title runs into the margin & a long path in a narrow column & a wider column (smaller code width) \\
			\hCode{\textbackslash hCodeJ} fails in a frame title & listings inside an argument & \hCode{\textbackslash hCode} \\
			hblibCodeoutput warns ``made for beamer'' & used in an article & use hblibCommon there \\
			\bottomrule
		\end{tabular}
\end{table}%
% +++++++++++++++++++++++++++++++++++++++ A tbl
\end{frame}
% ~~~~~~~~~~~~~~~~~~~~~~~~~~~~~~~~~~~~~~ A frm




% ~~~~~~~~~~~~~~~~~~~~~~~~~~~~~~~~~~~~~~ V frm
\begin{frame}[t,fragile] \frametitle{Quick reference}

% +++++++++++++++++++++++++++++++++++++++ V tbl
%: -Tbl. quick reference
\begin{table}
		\scriptsize
		\begin{tabular}{@{}L{0.62\textwidth}L{0.30\textwidth}@{}}
			\toprule
			\textbf{Command} & \textbf{Shows} \\
			\midrule
			\hCode{\textbackslash lstcode[code opts]\{code\}\{Language\}} & code \\
			\hCode{\textbackslash lstoutput[output opts]\{output\}} & output \\
			\hCode{\textbackslash lstterminal[opts]\{transcript\}} & a terminal session \\
			\hCode{\textbackslash lstcodedown[code opts][output opts]\{code\}\{Language\}\{output\}} & output below \\
			\hCode{\textbackslash lstcoderight[width][code opts][output opts]\{code\}\{Language\}\{output\}} & output right \\
			\hCode{\textbackslash lstcodedownstep[code opts][output opts]\{code\}\{Language\}\{output\}} & below, on slide 2 \\
			\hCode{\textbackslash lstcoderightstep[width][code opts][output opts]\{code\}\{Language\}\{output\}} & right, on slide 2 \\
			\hCode{\textbackslash lstcodeoverlaystep[code opts][output opts][width]\{x,y\}\{code\}\{Language\}\{output\}} & on top, slide 2 \\
			\hCode{\textbackslash slidegrid} & 1\,cm grid \\
			\hCode{\textbackslash hCodeJ\{\}}, \hCode{\textbackslash code\{\}}, \hCode{\textbackslash hCodeB\{\}} & inline Java, bash \\
			\hCode{\textbackslash hbLira}, \hCode{\textbackslash hbRupee} & \hbLira, \hbRupee \\
			styles: codestyle, outputstyle, terminalstyle, codesnippet, texsnippet & \hCode{style=\dots} \\
			\bottomrule
		\end{tabular}
\end{table}%
% +++++++++++++++++++++++++++++++++++++++ A tbl
\end{frame}
% ~~~~~~~~~~~~~~~~~~~~~~~~~~~~~~~~~~~~~~ A frm



\end{document}
